\documentclass[12pt]{article}
\usepackage[utf8]{inputenc}
\usepackage[russian]{babel}
\usepackage{amsmath}
\usepackage{amssymb}
\usepackage{geometry}
\usepackage{listings}
\usepackage{xcolor}

\geometry{a4paper, left=20mm, right=20mm, top=20mm, bottom=20mm}

% Настройка для кода Maple
\lstset{
    basicstyle=\ttfamily\small,
    keywordstyle=\color{blue},
    commentstyle=\color{green!60!black},
    stringstyle=\color{red},
    frame=single,
    breaklines=true,
    backgroundcolor=\color{gray!5}
}

\begin{document}

\section*{Решение}

\[
M = \{(x_1, x_2) \in \mathbb{R}^2 \mid x_1 \le 0,\ x_2 \le 0,\ x_1^3 + x_2^3 \le 1\}, \qquad x_0 = (0, 0).
\]

\subsection*{1. Ограничения и их активность в $x_0$}
\begin{align*}
g_1(x_1, x_2) &= x_1 \le 0 &\implies g_1(0, 0) &= 0 \quad (\text{активно}), \\
g_2(x_1, x_2) &= x_2 \le 0 &\implies g_2(0, 0) &= 0 \quad (\text{активно}), \\
g_3(x_1, x_2) &= x_1^3 + x_2^3 - 1 \le 0 &\implies g_3(0, 0) &= 0^3 + 0^3 - 1 = -1 < 0 \quad (\text{неактивно}).
\end{align*}
\[
\forall x \in M: \quad x_1 \le 0,\ x_2 \le 0 \implies x_1^3 + x_2^3 \le 0 < 1 \implies I(x_0) = \{1, 2\}.
\]

\subsection*{2. Градиенты и регулярность}
\begin{align*}
\nabla g_1(x_0) &= \left. \begin{pmatrix} \dfrac{\partial g_1}{\partial x_1} \\[8pt] \dfrac{\partial g_1}{\partial x_2} \end{pmatrix} \right|_{(0,0)} = \begin{pmatrix} 1 \\ 0 \end{pmatrix}, \qquad
\nabla g_2(x_0) = \left. \begin{pmatrix} \dfrac{\partial g_2}{\partial x_1} \\[8pt] \dfrac{\partial g_2}{\partial x_2} \end{pmatrix} \right|_{(0,0)} = \begin{pmatrix} 0 \\ 1 \end{pmatrix}.
\end{align*}
\[
\det \begin{pmatrix} 1 & 0 \\ 0 & 1 \end{pmatrix} = 1 \neq 0 \implies \text{условие регулярности выполнено.}
\]

\subsection*{3. Условие касательного конуса ($\nabla g_i(x_0)^T v \le 0,\ i \in I(x_0)$)}
\begin{align*}
\nabla g_1(0, 0)^T v &= \begin{pmatrix} 1 & 0 \end{pmatrix} \begin{pmatrix} v_1 \\ v_2 \end{pmatrix} = 1 \cdot v_1 + 0 \cdot v_2 = v_1 \le 0, \\[4pt]
\nabla g_2(0, 0)^T v &= \begin{pmatrix} 0 & 1 \end{pmatrix} \begin{pmatrix} v_1 \\ v_2 \end{pmatrix} = 0 \cdot v_1 + 1 \cdot v_2 = v_2 \le 0.
\end{align*}

\subsection*{Ответ}
\[
T_{(0,0)}M = \left\{ v = (v_1, v_2) \in \mathbb{R}^2 \;\middle|\; v_1 \le 0,\ v_2 \le 0 \right\}.
\]

\end{document}

\]

\end{document}
